%% Deliverable 2b — IEEE two-column conference paper
%% Compile with:  pdflatex Deliverable_2b_IEEE_Conference_Paper.tex   (twice, for the TOC)
%% Requires IEEEtran.cls (standard in every TeX Live / MiKTeX install).
%%
%% This file compiles to 12 pages in IEEE two-column format, references
%% included. It was 12 pages: the abstract was 500 words against IEEE's
%% 150--250, the results prose was a data dump, the related work carried
%% every prior-art finding twice over, and the bibliography spelled out every
%% author of every paper. It was then cut to 9 and to 10 as the related work
%% gained the 2026 literature that bears directly on the claims, and is now
%% 12 again because the paper gained a second corpus -- 110 United States
%% federal regulatory documents read at two past dates -- with the ablation
%% replicated on it, the conflict detector's precision measured as two
%% numbers rather than one, and a 2x2 factorial over the annotation rubric.
%% The number here is the measured one and the budget is recorded in
%% deliverables/scripts/check_latex.py with its reason.
%% The technical report (Deliverable 2a) and the journal version
%% (Deliverable 2c) carry the same material at their own lengths: this file
%% is the condensed venue version, and nothing was dropped from the thesis,
%% only from this one.
%%
%% The six figures live in deliverables/figures/ and are included by
%% relative path, so that directory has to travel with this file.

\documentclass[conference]{IEEEtran}

\usepackage{cite}
\usepackage{amsmath,amssymb,amsfonts}
\usepackage{algorithmic}
\usepackage{graphicx}
\usepackage{textcomp}
\usepackage{xcolor}
\usepackage{booktabs}
\usepackage{url}

\hyphenation{re-tri-ev-al}

\begin{document}

\title{An Evidence-Grounded Retrieval-Augmented Generation
Framework for Industrial Decision Support:
Architecture, Evaluation Protocol, and Failure-Mode Characterisation}

\author{\IEEEauthorblockN{Ritik Raj}
\IEEEauthorblockA{Department of Industrial and Systems Engineering,
Indian Institute of Technology Kharagpur, Kharagpur, India}
\and
\IEEEauthorblockN{Supervisor: Dr. Sri Krishna Kumar}
\IEEEauthorblockA{Department of Industrial and Systems Engineering,
Indian Institute of Technology Kharagpur, Kharagpur, India}}

\maketitle

\begin{abstract}
Industrial practitioners decide from paperwork: standard operating
procedures, failure-mode analyses and control plans. Retrieval-augmented
generation (RAG) grounds an answer in retrieved passages, but existing systems
trust the retriever unconditionally, ignore modality diversity, and silently
select one source when documents conflict. We present a five-mechanism
framework, every component of which has published prior art we position
against rather than claim priority over: a code-level, non-bypassable
retrieval-confidence gate that returns an explicit no-evidence response
without invoking the model; per-modality retrieval quotas; hybrid
lexical--dense retrieval fused by reciprocal rank fusion and reranked by a
cross-encoder; contradiction surfacing through a joint NLI and topical gate;
and per-claim three-band confidence tagging. On an encyclopaedic industrial
corpus, four properties of the original specification were wrong: a category
pre-filter on one retrieval leg is a no-op, a general-purpose sentence encoder
costs the fusion its ordering, the modality quota belongs before the
cross-encoder, and a reader confirms under two per cent of the contradiction
gate's alarms. We then repeat every measurement on a second corpus of 110
United States federal regulations read at two past dates, which supplies
conflicts whose types are known by construction. There a conflict detector's
precision splits into two numbers nine-fold apart, so the mechanism finds the
conflicts and reports them as the wrong kind; and a $2\times2$ factorial shows
the annotation rubric, not the reader's specialisation, moved that number by a
factor of nine. Auditing our own benchmark showed its unanswerable class
separable from the question text alone, so we rebuilt it; the gate now refuses
{{UNANS_HARD_REFUSED_PCT}} of the repaired hard class against the
{{UNANS_BEFORE_REFUSED_PCT}} the defective class suggested.
\end{abstract}

\begin{IEEEkeywords}
Retrieval-augmented generation, industrial decision support, hallucination
reduction, confidence gating, selective prediction, hybrid retrieval,
per-claim attribution, evaluation methodology.
\end{IEEEkeywords}

\section{Introduction}

Industrial decision-making depends on timely access to evidence in technical
documentation. Practitioners across quality, maintenance, logistics and
process operations consult standards such as ISO 9001, GS1 specifications,
Incoterms rules and original-equipment-manufacturer service procedures. It
has grown with industrial digitalisation faster than it can be synthesised,
so decisions are made with partial information, at a cost to efficiency,
compliance and safety.

The proposal behind this work described a twelve-category industrial benchmark
spanning industrial engineering, machine learning and data analytics. This
paper evaluates six of the nine industrial-engineering categories, on three
corpora that could be obtained and redistributed, and claims nothing about
the twelve-category benchmark: the three machine-learning categories are
numerical series-reasoning problems rather than document-retrieval problems,
and three industrial categories are descoped because the binding artefacts the
proposal names are proprietary and not redistributable. The reconciliation, with
the reason for every row, is in the supplementary documentation.

Large language models (LLMs) and retrieval-augmented generation (RAG) have
been proposed to address the cost of locating, cross-referencing and
verifying evidence under time pressure. Three properties industrial users
require remain absent from generic systems: no \emph{evidence trail} saying
which passage supports which claim; no \emph{confidence check}, so the system
answers when retrieval is weak; and no \emph{contradiction handling}, so when
two accepted passages disagree it silently selects one.

SQuAD 2.0 established the general form of the problem with an evaluation
requiring a system to answer answerable questions correctly \emph{and}
decline unanswerable ones, and reported a strong neural system at 86\% F1 on
SQuAD but only 66\% on SQuAD 2.0. Twenty points are lost to confidently
answering unanswerable questions in a general-domain setting; no such loss
is acceptable for a plant's reorder policy.

We measure which reliability mechanisms pay for themselves in industrial
text; the measurement, not the architecture, is the contribution, in three
parts. The \emph{architectural} part is incremental: by late 2026 prior art
publishes a code-level pre-generation gate with a calibrated decision policy
(Sieve and Sage~\cite{won2026sievesage};
HALT-RAG~\cite{goswami2025haltrag}), set-level sufficiency verification and
structural certificates (SURE-RAG~\cite{qiu2026surerag};
Certify~\cite{zhang2026certify}), and the argument that modality routing is
performed at the wrong decision point (Modality Relevance is not Modality
Utility~\cite{chen2026modalityutility}), so we position against it, not
claiming priority. The
\emph{empirical} part is a failure-mode characterisation with negative
results: the dense leg costs the fusion its ordering until the encoder is
replaced, the modality quota's pipeline position inverts the framework's
design rule, and the confidence gate cannot separate absence of evidence
from retrieval failure, making it an answer-in-context
detector~\cite{li2026relevance}, not a sufficiency judge. The
\emph{methodological} part is an evaluation protocol
treating the gate as a selective-prediction policy, reporting a
risk--coverage curve and a false-positive rate at a stated operating point
against two degenerate baselines (never refuse, always refuse) rather than
the sweep endpoints.

A fourth part is a metric contribution, and it is the paper's clearest new
result. A conflict detector's precision, reported as one number, conflates
two questions: \emph{is the alarm pointing at something real}, and \emph{is
the branch that raised it the branch that owns the conflict}. We measure both,
on a second corpus of United States federal regulations read at two past
dates, where the same provision at two issues \emph{is} a conflict with its
type known by construction and so needs no annotation. The two numbers come
apart by a factor of nine, and the mechanism finds the conflicts and reports
them as the wrong kind.

A second methodological contribution follows from auditing the benchmark:
its unanswerable class proved separable from the question text alone, making
the headline refusal rate an upper bound rather than a measurement, so the
class was rebuilt to make unanswerability decidable only from the evidence
(Section~\ref{sec:unansaudit}). The repair changes what the headline number
means, so it is a contribution rather than a caveat.

\section{Related Work}

\subsection{Retrieval-augmented generation and its failure modes}

RAG~\cite{lewis2020rag} set the retrieve-then-generate paradigm; dense passage
retrieval~\cite{karpukhin2020dpr} set dual-encoder search, beating a strong
Lucene-BM25 baseline by 9--19 percentage points absolute in top-20 passage
retrieval accuracy. Self-RAG~\cite{asai2024selfrag} trains reflection tokens
deciding when to retrieve and whether a passage supports the output;
CRAG~\cite{yan2024crag} inserts a retrieval evaluator whose confidence degree
triggers one of four corrective actions; FLARE~\cite{jiang2023flare} and
DRAGIN~\cite{su2024dragin} trigger retrieval mid-generation from the model's
own information need; Adaptive-RAG~\cite{jeong2024adaptive} routes queries by
complexity. Our gate is not \emph{when to retrieve} but \emph{whether to
answer}, and no by default when the evidence is weak. The premise is
published: enterprise RAG needs purpose-built architectures and its own
evaluation framework, and three of our five mechanisms have been implemented
independently in cross-modal financial QA. Neither the premise nor the
mechanisms are claimed as new.

\subsection{Confidence, calibration and abstention}

SQuAD 2.0~\cite{rajpurkar2018squad2} set the unanswerable class; Kadath et
al.~\cite{kadavath2022know} showed verbalised confidence is informative;
semantic entropy~\cite{kuhn2023semantic} and Farquhar et
al.~\cite{farquhar2024semantic} separate \emph{confabulation} from
\emph{consistent error}. Joren et al.~\cite{joren2024sufficient} is the
closest result: large models answer well on sufficient context but answer
wrongly rather than abstain when the context is insufficient, so a gate on the
\emph{generated answer} inherits that failure while a gate on the
\emph{retrieval} attacks the right variable. Decian et
al.~\cite{decian2026dynamic} combine hybrid retrieval with an entropy-based
dynamic threshold; Gei{\ss}ler et al.~\cite{geissler2026dependable} obtain
answer-quality gains up to 6\% and inconsistency detection up to 77\%, with
guarantees conditional on exchangeability, which the same machinery extends to
pre-generation filtering~\cite{chen2026principled} and per-claim
factuality~\cite{feng2025conformal,claimconformal2026}; entropy-based
uncertainty also has a model-dependent failure mode~\cite{phillips2026entropy}.
Thakur et al.~\cite{thakur2026declinecraft} show query rewriting and reranking
\emph{lower} the correct-abstention rate for all 12 models. SYNCHECK
~\cite{syncheck2024} is the sharpest challenge to D1 --- decoding beats
abstention by over 10\% on faithfulness --- but it asks whether a generated
sentence is faithful to the context, while the gate asks whether the context
contains enough to answer at all, settled before generation, since a sentence
from inadequate evidence cannot be repaired afterwards. Long-context
architectures are ahead of semantic RAG in manufacturing safety at a materially
higher token cost~\cite{longcontext2026tokentax}; published defences against
query-time poisoning can measure benchmark construction rather than
filtering~\cite{ranking2026notdeciding}; better retrieval reduces abstention,
so the operating point must be re-fitted per configuration. I-CALM
~\cite{zong2026icalm} recovers each abstention's counterfactual answer, with
coverage, error-capture and correct-retention rates at matched coverage,
although coverage often collapses; TAPS~\cite{zhang2026taps} splits abstention
risk into answering and distinguishing ability; UAEval4RAG~\cite{peng2025uaeval4rag}
treats prompt restrictiveness as an explicit factor, which we follow by
comparing arms at matched coverage. The Selection Effect
~\cite{tikhomirov2026selection} shows encoder gating can forward a slice where
the downstream component is no better than the baseline, and
Recognition--Refusal Misalignment~\cite{du2026recognition} shows recognising
impossibility is nearly orthogonal to routing it to abstention, so our null
result is reported as the prompt having already realised the gate's function on
that slice. Models overcommit and rarely abstain~\cite{cocchieri2026overcommit},
so one generator cannot carry a general claim. Jev-as-a-Judge
~\cite{rosario2026jev} is statistically indistinguishable on the conventional
metric yet disagrees on 15.2\% of claims and admits very different unsupported
rates, 23.2\% against 14.1\%, on the same 495 human-labelled claims; a matched
random-escalation control does as well, so we run that control too: an arm
refusing as many questions as the fitted gate, picked at random. RCL
~\cite{ravuri2026rcl}, HIRA~\cite{tian2026hira} and CiteGuard-RAG
~\cite{zheng2026citeguard} place a screening layer between retrieval and
generation, so the mechanism is prior art; what we claim is the
\emph{measurement} --- per-claim confidence banding with a cited evidence
trail, on industrial text, at the level of the answer claim.

Several 2026 papers pre-empt the gate's architecture, and this work is
positioned against them. Sieve and Sage~\cite{won2026sievesage} screen the
retrieved set in a module outside the model, the same pre-LLM placement used
here; SURE-RAG~\cite{qiu2026surerag} makes sufficiency a set-level property
that independent passage scoring cannot detect; Certify~\cite{zhang2026certify}
makes it structural, with Horn-closure certificates reported to beat learned
sufficiency judges; and other work computes confidence in the orchestrator,
never asking the model. What is left distinct is a calibrated detector with a
stated operating point and a risk--coverage curve, not a certificate. The label
follows the same literature, which reserves \emph{sufficiency} for the
set-level property~\cite{li2026relevance}, so the quantity predicted here is
named \emph{answer-in-context detectability}; and since a model's apparent
willingness to abstain can be produced by the prompt rather than by genuine
uncertainty~\cite{ling2026promptartifact}, the gate's contribution is measured
against a prompt-conditioned baseline under two prompt conditions differing by
a single sentence.

\subsection{Multimodal retrieval}

ColPali~\cite{faysse2025colpali} introduced page-image retrieval with
multi-vector vision-language embeddings and late-interaction matching.
VisDoMRAG~\cite{suri2024visdom} fuses modalities under a consistency
constraint, our closest prior art on modality balancing.
T\textsuperscript{2}-RAGBench~\cite{strich2026t2} provides a text-and-table
benchmark of 23{,}088 expert-validated triples, hybrid BM25 best on
text-plus-table queries. Modality-balanced retrieval is already in the
literature, so our claim concerns the \emph{mechanism}, a hard non-bypassable
quota on the final context, and the domain. \emph{Modality Relevance is not
Modality Utility}~\cite{chen2026modalityutility} shows modality relevance a
weak predictor of modality utility, so a modality quota is mis-specified unless
utility is ranked; a placebo-controlled ablation of four retrieval mechanisms
across two corpora \cite{sharma2026placebo} finds naive hierarchical and
structure-aware variants can hurt, corpus-conditionally. The trade has a
published mechanism: constraint-induced displacement shows
that candidates conflicting with the query's period, segment, metric scope or
table scope consume top-$k$ slots and displace the evidence that would have
answered it, costing Recall@10 by 0.147 against an equally difficult
compatible pool. The quota studied here displaces in the same way, so the
measurement below is of a published mechanism.

\paragraph{Lexical or dense on technical and regulatory text.} That lexical
retrieval beats an off-the-shelf dense encoder on legal and regulatory text
is established, and this paper says so in the first line of the section that
reports it. On Court of Justice of the European Union passages, BM25
``surpassing off-the-shelf dense models in 4 out of 7 performance metrics'',
and the same study reports that fine-tuning a dense model on domain-specific
data overtakes BM25~\cite{mori2025lexical}; statutory article retrieval in
Dutch and French finds BM25 ``remains a competitive baseline'' against
zero-shot dense models~\cite{lotfi2024bsard}; and a benchmark of ten
retrieval strategies over 23{,}088 financial text-and-table queries reports
that BM25 ``outperforms state-of-the-art dense retrieval on financial
documents, challenging the common assumption that semantic search universally
dominat[es]''~\cite{akarsu2026bm25}. The counter-case is reported too: on
Greek statutory retrieval, vanilla dense retrieval outperforms vanilla sparse
at Recall@100, though query reformulation closes the gap~\cite{beta2026greekbar}.
TeleEmbedBench~\cite{gajjar2026teleembed} exists because general-purpose
benchmarks fail to capture the dense, acronym-heavy nature of technical
corpora. Three consequences follow: the comparison is \textbf{not} claimed as
new and does not appear in the abstract; the dense leg here is a
general-purpose encoder a domain-adapted one could overtake, so the result is
a boundary condition on this pipeline rather than a claim about dense
retrieval as a class; and the \emph{corollary} is the new part --- on this
corpus the lexical leg carried the pipeline, and neither the fusion nor the
category pre-filter added anything measurable on top of it.

\subsection{Attribution, evaluation, and the failure of ``reported a number''}

ALCE~\cite{gao2023alce} defined citation recall and precision, with even the
best models lacking complete citation support for ELI5 half the time;
RAGAS~\cite{es2024ragas} defined faithfulness, answer relevancy and context
precision; ARES~\cite{saadfalcon2024ares} provides a reference-free framework;
FActScore~\cite{min2023factscore} decomposes generations into atomic facts,
with ChatGPT biographies only 58\% factually supported. Two 2026 papers change
how these are reported: benchmark construction cues can dominate measured
filter performance in query-time RAG poisoning defences, AUROC falling from
0.953 to 0.758 once a cue is removed; and cheap LLM judges give AUROC as low as
0.603, with position bias flipping 66 of 80 queries. We adopt both as
requirements on our own evaluation. A third finding constrains the labels
rather than the metrics: automated labelers tend to conflate reference
faithfulness with factual correctness~\cite{valjakka2026labeling}, which is
why the labels here are audited rather than trusted. Nor is this the first
industrial benchmark: comparable corpora now cover manufacturing safety,
automotive requirements engineering, industrial standards, semiconductor QA,
regulated documents and data-lake QA, with RealDocBench and IndusSpec-QA the
closest in kind~\cite{joshi2026realdocbench,park2026induspecqa}. Those are
expert-annotated and domain-authentic where this one is derived and
encyclopaedic; the contribution here is the refusal half.

A fourth finding constrains what a model-annotated measurement may be
reported to mean, and it bites hardest in the second corpus this paper
introduces. The annotation guideline is a \emph{first-class error source} in
model annotation, not a formality: SILICON decomposes LLM-annotation
measurement error into four sources, the first being guideline-induced error
from inconsistent annotation criteria, and frames model identity as an error
to control rather than a capability to buy~\cite{cheng2024silicon}; a 2026
ACL main-conference paper reports that models ``often struggle with the
specialized conventions of gold-standard benchmarks'' and proposes guideline
refinement as an alignment mechanism~\cite{kim2026guidelines}. The principle
is published; what this paper adds is a factor of nine, and that it lands on
the \emph{reported system measurement} rather than only on the labels.

\subsection{Contradiction and conflict detection}

NLI is the standard instrument~\cite{williams2018mnli,nie2020anli}.
ContextConflict~\cite{yang2026contextconflict}: six conflict types, factual,
inferential, temporal, granularity, perspective and ambiguity, of which only
two are genuine contradictions, the strongest published justification for a
\emph{topical} gate and a conflict \emph{type}, not a binary flag.
ConflictScore~\cite{liu2026conflictscore}: decomposes responses into atomic
claims and labels each against each grounding document.

\paragraph{What changed in 2026, and what it forces us to narrow.} GRAB-RAG
\cite{setiawan2026grabrag} finds that models still answer \emph{41.6\%} of
misleading questions under explicit abstention prompting, and concludes that
prompt-based abstention asks whether context is sufficient rather than
whether it is correct. This work reproduces that finding in industrial text,
so the negative result corroborates it in a new domain rather than surprising
locally; with SURE-RAG, Certify and Sieve and Sage as above, our gate is
consequently an \emph{answer-in-context detector}, not a sufficiency judge.
HALT-RAG \cite{goswami2025haltrag}, the closest published pre-emption,
combines a calibrated NLI ensemble with lexical signals under a
precision-constrained abstention policy, and Kontrast \cite{wei2026kontrast}
formalises the temporal-change distinction we implement as a supersession
check; HALT-RAG is a \emph{post-hoc} verifier that never retrieves, and its
Future Work names integrating detectors like it in a live RAG pipeline as the
critical next step, which is our Section~\ref{sec:industrial}, where the
detector is measured end-to-end on a corpus that contains conflicts.
Beyond the Reranker
\cite{singh2026beyondreranker} finds that retrieval enhancements stop paying
once a strong reranker is present, and that the gains come almost exclusively
from homogeneous, predominantly Wikipedia prose corpora, which ours is not.
WIKICOLLIDE, with CLAIRE \cite{semnani2025wikicollide}, measures at least
3.3\% of English Wikipedia facts as contradicting another fact, and the
inconsistency propagates into 7.3\% of FEVEROUS examples, a corpus built from
industrial documentation; theirs is a corpus-level measurement whose unit is a
pair of facts anywhere in the corpus, ours a surface mechanism over the
passages retrieved for one query, so no corpus-level consistency rate is
claimed here and the absence of contradictions in this corpus is a property of
a narrow slice of encyclopaedic prose rather than a finding about Wikipedia.

\paragraph{Point-in-time regulatory corpora, and what they are for.}
Building a corpus that provably contains conflicts, with their types known
without annotation, is a published method. NitiBench+\cite{akarajaradwong2026nitibench}
constructs a temporal legal question-answering benchmark across 17 Thai
domains from point-in-time statute snapshots; Asking For An Old Friend
\cite{prior2026oldfriend} diagnoses temporal failure modes in statutory
question answering over the same kind of snapshot; and When Do LLMs Apply the
Wrong Law?\cite{huang2026wronglaw} applies versioned statutes to the question
of which law applies. This paper does not claim the method. Those works use a
versioned corpus to decide \emph{which law applies}; here the same provision
read at two issues \emph{is} a conflict, with its type known by construction,
so the published method is pointed at a new task. On the detector side,
\emph{Detecting Is Not Resolving}~\cite{yu2026detecting} names the gap this
paper measures --- a system can detect that two sources conflict and still
fail to say which one governs --- and DRAGged into Conflicts
\cite{cattan2025dragged} builds a benchmark of conflicting sources and
addresses them with an intervention rather than a taxonomy.

\section{The Framework}

The framework has three stages: ingestion (offline), retrieval and generation
(online).

\begin{figure}[t]
\centering
\begin{minipage}{0.92\linewidth}
\small
\textbf{Stage 1 --- Ingestion (offline).}
parse $\rightarrow$ chunk (text\,$\vert$\,table\,$\vert$\,chart) $\rightarrow$
tag category $\rightarrow$ embed $\rightarrow$ index.

\textbf{Stage 2 --- Retrieval (online, per query).}
form check $\rightarrow$ embed query $\rightarrow$ category pre-filter
$\rightarrow$ BM25 leg $+$ dense
leg $\rightarrow$ reciprocal rank fusion $\rightarrow$ cross-encoder rerank
$\rightarrow$ \textbf{confidence gate} $\rightarrow$ modality quota.

\textbf{Stage 3 --- Generation (online, per accepted query).}
per-claim tagging $\rightarrow$ citations $\rightarrow$ contradiction flag
$\rightarrow$ tagged answer.
\end{minipage}
\vspace{2pt}
\caption{The three-stage pipeline. The confidence gate is the load-bearing
element: when it rejects, the language model is never invoked.}
\label{fig:pipeline}
\end{figure}

\subsection{Ingestion and retrieval}

A PDF is a visual object: a page lays out text blocks, ruled tables and
figures, and becomes a table only once the ruling lines are understood. An
industrial corpus keeps its binding requirements in tables, so the parser is
structural: blocks with bounding boxes are classified as \emph{text},
\emph{table} or \emph{chart}. Tables are serialised to markdown, merged cells
forward-filled, never split; text is chunked into 120-word windows with
30-word overlap. Modality composition is measured and reported, since it
predicts whether modality-balancing can have any effect.

The lexical leg is BM25~\cite{robertson2004prf}, implemented from first
principles:
\begin{equation}
\mathrm{BM25}(q,D)=\sum_{t\in Q}\mathrm{IDF}(t)\,
\frac{f(t,D)(k_1+1)}{f(t,D)+k_1\big(1-b+b|D|/\overline{dl}\big)},
\label{eq:bm25}
\end{equation}
with $k_1=1.2$ and $b=0.75$. The dense leg embeds the query and every chunk
with {{EMBED_MODEL_TEX}} ({{EMBED_DIM}} dimensions, L2-normalised) and ranks by
cosine similarity. The two rankings are fused by Reciprocal Rank
Fusion~\cite{cormack2009rrf},
\begin{equation}
\mathrm{RRF}(d)=\sum_i \frac{1}{k+r_i(d)},\qquad k=60,
\label{eq:rrf}
\end{equation}
because BM25 scores, cosine similarities and cross-encoder logits are
incomparable: summing them needs a fitted weighting that does not transfer
across corpora, whereas rank is scale-free. Bruch et
al.~\cite{bruch2023fusion} found RRF at least as good as learned
score-combination across 23 runs, with \emph{tuned} weights failing to
generalise out of domain. A cross-encoder then reranks the shortlist; its
output is a \emph{logit, not a probability}, whose distribution shifts with
query difficulty.

\subsection{The confidence gate}

With $R$ the event that a chunk is relevant and $d$ the evidence, Bayes'
theorem in odds form is
\begin{equation}
P(R\mid d)=\frac{P(R)}{P(d\mid\neg R)}\cdot\frac{P(d\mid R)}{P(d\mid\neg R)},
\label{eq:bayes}
\end{equation}
and the gate uses a logistic model on log-odds,
\begin{equation}
P(R\mid d)=\sigma(\beta_0+\beta_1 x),\qquad \sigma(z)=\frac{1}{1+e^{-z}},
\label{eq:posterior}
\end{equation}
fitted by maximum likelihood on a split held out by \emph{query}, not by
candidate.

The gate is a \emph{query-level} decision: not whether this candidate is the
gold chunk, but whether the context contains the answer at all. With about
fifty candidates per query and one gold chunk, chunk-level positives have a
base rate of roughly 2\%, so a classifier answering ``nothing is relevant''
is right 98\% of the time; fitting against that label yields a large negative
intercept and a nearly flat slope, so calibration metrics look excellent:
they measure the base rate, not the filter. The label is therefore 1 when the
context contains a gold chunk, found or missed: in both cases the correct
behaviour is to refuse. The five features, all available before generation:
the top cross-encoder score, the mean of the top three scores, the best BM25
and RRF scores, and the number of candidates clearing a low relevance bar.

The gate is control flow:
\begin{equation}
\begin{aligned}
&\text{if } P(\text{context sufficient}) < \tau:\ \text{return no-evidence response},\\[-1pt]
&\text{else } \mathcal{C}=\text{accepted chunks}.
\end{aligned}
\label{eq:gate}
\end{equation}
It cannot be talked past: no output reaches it; its refusal is auditable
(``maximum calibrated relevance probability $0.312$ against threshold
$0.500$'' is checkable, ``the model seemed unsure'' is not); and it is
portable without retraining, which matters across industrial categories:
Self-RAG's reflection tokens would be re-learned per category, and Huang
et al.~\cite{huang2024selfcorrect} show self-correction degrades without
external feedback. The \emph{calibration} is a separate matter and is not
portable: the threshold is a fitted parameter, and it is re-fitted per corpus
and per retriever, which is measured rather than assumed. CRAG, the nearest prior art, responds to low confidence by
improving the evidence and continuing, and calls the LLM in every branch.

\subsection{The other three mechanisms}

A form check runs before retrieval, refusing an empty query, a query of only
function words or deictics, and a query ending on a word that cannot end a
question; it is measured rather than asserted, and on
{{FORMVAL_CORPUS}} it refused {{FORMVAL_STEM_MAL_REJECTED}} of
{{FORMVAL_STEM_MAL_N}} hand-labelled malformed stems
({{FORMVAL_STEM_MAL_RATE}}) against {{FORMVAL_STEM_GOOD_REJECTED}} false
rejections among {{FORMVAL_STEM_GOOD_N}} well-formed stems.

The \emph{modality quota} is a floor on the share of the context taken by
tables (0.15) and charts (0.05), implemented as a swap. Its
\emph{pipeline position} is decisive: before the reranker it proposes a
membership for the final ten, which the reranker can reorder but not change,
so the floor binds; after the reranker it swaps the lowest-ranked chunk for
one the reranker already ranked below the cut, which tends to import the
wrong chunk and displace a good one. Section~V reports the measured
difference: {{QUOTA_PLACEMENT_CLAIM}}, and
{{QUOTA_PLACEMENT_SIGNIFICANCE}} The
shipped system applies the quota {{SHIPPED_QUOTA_PLACEMENT}} --- it was built
the other way round and was moved on 16 November 2026, after which
{{PLACEMENT_SCOPE}}.

Two accepted chunks conflict only when two gates agree,
\begin{equation}
\mathrm{flag}(A,B)=\mathrm{contradiction}(A,B)\wedge\mathrm{same\_subject}(A,B),
\label{eq:flag}
\end{equation}
where \texttt{contradiction} is the NLI probability from a DeBERTa-v3 MNLI
model and \texttt{same\_subject} combines embedding cosine similarity
(threshold 0.55) with Jaccard overlap of concrete identifiers (threshold
0.10). The conjunction lets the two error types be traded off separately. A
\emph{numeric conflict}, two passages attaching different values to the same
quantity, is an explicit contradiction regardless of confidence. A
\emph{revision-supersession check} is required: ``Rev 3 sets EOQ at 187 units''
and ``Rev 4 sets EOQ at 95 units'' are both true (Rev 4 replaced Rev 3), so
calling that a contradiction is wrong.

Answers are decomposed into atomic claims by sentence and split at discourse
markers. For each claim the NLI model gives an entailment probability $e_i$
per retrieved chunk, combined by the noisy-OR model: each chunk is a fallible
witness and the claim is supported if one covers it,
\begin{equation}
q = 1-\prod_i (1-e_i).
\label{eq:noisyor}
\end{equation}
The band is cut on $q$: $q\geq0.80$ is \emph{grounded}, $q\geq0.45$
\emph{inferred}, otherwise \emph{uncertain}; information gain in bits is
reported beside it, and is symmetric about $q=0.5$.

\section{Evaluation Protocol}
\label{sec:protocol}

\paragraph{Benchmark.} We evaluate on a \emph{silver-standard} benchmark
derived from an independently constructed public industrial corpus, not the
300-query expert-annotated SupplyChainQA v7 of the earlier project report,
which is proprietary and unreleased. Labels are derived from the document: a
table cell becomes a lookup question with the table chunk as gold, a
definitional or numeric sentence a question with its chunk as gold. This is
slightly optimistic, since the derived gold chunk always contains the answer
somewhere, inflating \emph{absolute} recall without biasing the
\emph{comparison} between systems scored against the same labels. The
benchmark contains deliberately unanswerable questions on topics absent from
the corpus, without which no confidence-gating experiment is meaningful: a
system answering everything scores 100\% recall and 0\% hallucination.

\paragraph{The corpus and its provenance.} Every measurement here uses a public
corpus of {{N_DOCUMENTS}} industrial-engineering articles
({{N_CHUNKS_TEX}} chunks, {{N_CATEGORIES_WORD}} categories), assembled from
English Wikipedia via the MediaWiki API, licensed CC BY-SA 4.0. It
substitutes for the specification's corpus of 634 chunks across 24
proprietary supply-chain documents, which cannot be redistributed and so
cannot ship with the paper or be re-run. The substitute ships in full with a
manifest giving each document's Wikipedia URL, upstream revision, the SHA-256
of the bytes used, and its chunk count. The focus is one industrial domain,
supply chain management ({{N_FOCUS_DOCUMENTS}} documents,
{{N_FOCUS_CHUNKS_TEX}} chunks, {{N_FOCUS_QUERIES}} questions); the rest
support the cross-category transfer objective.

\paragraph{What ``silver-standard'' means here.} The first version was worse
than that word suggests: a quarter of its answerable questions were malformed,
because the generator's subject pattern could match mid-clause or after a
truncated word, giving questions like
``What value is stated for a teep in the source documentation?''. A benchmark
whose questions are a quarter malformed measures nothing. The generator now
anchors each match to a sentence start, rejects subjects with a verb, pronoun
or dangling preposition, rejects table cells with markup, and keeps acronyms
in their own case, so questions ask about TEEP rather than ``teep''. The result
is a smaller benchmark, and it remains silver-standard: labels are derived,
not annotated.

\paragraph{Metrics and statistics.} Retrieval is measured by Recall@$k$,
Precision@$k$, MRR, nDCG@$k$ with graded relevance
(3 for a chunk carrying a binding specification, 1 otherwise), and
\emph{per-modality recall}; generation by faithfulness, citation precision
and citation recall. Queries with no gold set are excluded from recall and
nDCG aggregates, not averaged in as zero. Comparisons are paired on the same
queries: binary outcomes by McNemar's exact test, continuous per-query
outcomes by the Wilcoxon signed-rank test, which suits bounded, left-skewed
faithfulness. Effect sizes (Cohen's $d_z$, Cliff's $\delta$) accompany every
$p$-value, and bootstrap percentile intervals cover the headline figure.

\paragraph{Gate as a selective-prediction policy.} We sweep $\tau$ and report
the fraction of answerable questions answered, the fraction of unanswerable
questions refused, abstention precision and coverage. The operating point
follows a stated rule, the highest $\tau$ still answering
at least 90\% of answerable questions, not a favourable metric, and the full
risk--coverage curve is reported. Because reranking in hybrid retrieval
lowers the correct abstention rate, the calibration is fitted
\emph{separately for each retrieval configuration compared}.

\section{Results}

\subsection{Retrieval ablation}

Seven systems were compared on the same {{N_FOCUS_QUERIES}} queries over the
{{N_FOCUS_CHUNKS_TEX}}-chunk, {{N_FOCUS_DOCUMENTS}}-document focus corpus,
supply chain management. The six-category corpus carries the transfer
experiment, which needs more than one category. Table~\ref{tab:ablation} reports
the results; paired Wilcoxon tests between adjacent rungs appear in the last
column.

\begin{table}[t]
\caption{Retrieval ablation on {{N_FOCUS_QUERIES}} silver-standard queries over the
{{N_FOCUS_CHUNKS_TEX}}-chunk focus-domain corpus. $\Delta$R@10 is the paired Wilcoxon signed-rank
difference to the preceding rung.}
\label{tab:ablation}
\centering
\small
\resizebox{\columnwidth}{!}{%
\begin{tabular}{lrrrrrr}
\toprule
System & R@5 & R@10 & R@20 & MRR & nDCG@10 & Tbl R@10 \\
% The rows between \midrule and \bottomrule are written by
% deliverables/scripts/make_tex_results.py from the measured ablation JSON.
% A hand-written copy of a generated table is a copy that goes stale
% silently, and this one had: it still carried the pre-encoder-swap numbers.
% The empty row is a placeholder so that the table compiles before the
% generator runs; it carries no measurement and is replaced either way.
\midrule
 & & & & & & \\
\bottomrule
\end{tabular}
\vspace{2pt}
\begin{tabular}{lrl}
\toprule
Rung & $\Delta$R@10 & $p$ (Wilcoxon) \\
\midrule
 & & \\
\bottomrule
\end{tabular}}
\end{table}

% BEGIN RESULTS-NUMERIC
% The block between these markers is generated from the experiment JSON by
% deliverables/scripts/make_tex_prose.py. Nothing is written here on purpose:
% a hand-written copy of a generated block is a copy that goes stale
% silently -- and it did, once, with numbers that no longer matched any run.
% END RESULTS-NUMERIC

\subsection{The confidence gate, contradiction surfacing and transfer}
The gate was evaluated as a selective-prediction policy over a threshold
sweep with the calibration fitted on a query-level 60/40 split. The fitted
coefficients differ substantially from the earlier project report's values
($\beta_0=2.943$, $\beta_1=1.680$, fitted on 30 queries of a different
corpus), which is not a contradiction but a demonstration that \emph{the
calibration is a fitted parameter of the deployment, not a constant of the
architecture}. The operating point trades coverage for correctness and both
sides must be reported: the always-answer baseline refuses no unanswerable
questions, so reporting only the refusal rate overstates the gain while
reporting only the answerable answer rate hides the cost.

The joint contradiction gate was compared against the NLI-only gate over every
pair of accepted chunks, with the prevalence of true conflicts and the
always-flag and never-flag baselines reported alongside. The same-subject check
suppresses precisely the pairs an NLI model labels \emph{contradiction} when
forced to produce a label on two passages about different sub-topics. Its
precision was then measured against reader labels rather than against a proxy.
{{CONTRA_PRECISION_CLAIM_SHORT}} {{CONTRA_PRECISION_MECHANISM_SHORT}} The
gate's findings are surfaced only when the measured precision of the branch
that raised them clears a calibration floor --- the same discipline the answer
gate applies to an answer --- which on this corpus leaves
{{CONTRA_SURFACED}} findings surfaced and {{CONTRA_WITHHELD}} withheld.
{{CONTRA_CONTROL_CLAIM}} The
revision-supersession check suppresses a further class whose importance is not
a matter of degree: telling a practitioner that a policy is inconsistent when
it was in fact updated is the single most damaging error available to the
system.

The gate's coefficients and threshold were fitted on one half of the six
industrial categories and applied unchanged to the held-out half, with a
leave-one-category-out variant. The \emph{ranking} of configurations transfers
across industrial sub-domains; the \emph{threshold} does not, and the error is
systematic --- a threshold fitted on one configuration is too permissive where it
is transferred, and too permissive in the direction that flatters the
configuration it is transferred to. At the tuning half's threshold applied
unchanged to the held-out half, the gate answers
{{XFER_HELD_ANSWERED}} of the answerable questions against
{{XFER_HELD_REFIT_ANSWERED}} when the weights are refitted there, and refuses
{{XFER_HELD_REFUSED}} of the unanswerable ones against
{{XFER_HELD_REFIT_REFUSED}}. This is the practical form of the decline in
abstention reported by Thakur et al., and it is why the calibration is
re-fitted per arm.

\subsection{A defect in our own benchmark, found by auditing it and then
repaired}
\label{sec:unansaudit}

The unanswerable half of a refusal benchmark is where the benchmark lives or
dies, and it was the half we audited last. The audit found a serious defect.
A classifier that saw only the question --- no retrieval, no evidence, no
reasoning about the corpus --- told the answerable questions from the
unanswerable ones at an ROC AUC of {{UNANS_BEFORE_AUC}}. The mechanism was
concrete: the half's designated hard case asked about an identifier verified
absent from the corpus (\texttt{SKU AST-5454}), and every one of those
questions carried an identifier shape that no answerable question carried, so
one regular expression separated the two classes perfectly, catching
{{UNANS_SHORTCUT_BEFORE}} of the class with no false positives at all. If a
system's unanswerability is visible in the question, a system refusing on the
question's shape scores as well as one reading the passages, and the benchmark
cannot tell them apart.

We rebuilt the half rather than reporting around it. The repaired class uses
the answerable class's own three question stems, its own source documents and
its own label machinery --- the definition branch finds a definition, the
numeric branch finds a value, a table cell holds the answer --- and is labelled
by the complement of exactly those tests. Each repaired question was further
drawn to have a \emph{structural twin} in the answerable class: the same stem,
the same subject size, the same shape markers, and no word in common;
{{UNANS_MATCH_RATE}} of them have one. Re-measured, the shape-only classifier
scores {{UNANS_SEP_AUC}} against a permutation null of {{UNANS_SEP_NULL}}, and
the one-rule shortcut carries {{UNANS_SHORTCUT_AFTER}}. A question-only policy
--- a model given the question and no documents, asked whether the corpus
could answer it --- discriminates the two classes at a ROC AUC of
{{UNANS_POLICY_BEFORE_AUC}} before the repair and {{UNANS_POLICY_AUC}} after,
where 0.5 is no discrimination at all. The same artefact has been reported in
the same task for a clinical question-answering benchmark, where the
unanswerable questions ``can often be discerned simply by filtering with
specific N-gram patterns'' and the evaluation is thereby compromised
\cite{yang2024ehrsql}; the shortcut-removal construction used here is the one
applied to insufficient-evidence examples elsewhere
\cite{mondal2026beforeanswering}; and framing a benchmark repair as the
contribution follows the precedent of \cite{brunello2026folio}. What is not
reported anywhere is the question-only separation of an answerable/unanswerable
split in a RAG \emph{refusal} benchmark, its effect on the reported refusal
rate, and a repair by regenerating the unanswerable class with answerable-style
templates \cite{gautam2023unanswerable}.

The refusal rate is now reported by stratum. At $\tau = {{TAU}}$ the gate
refuses {{UNANS_HARD_REFUSED_PCT}} of the corpus-anchored questions
({{UNANS_HARD_REFUSED_N}}) --- questions the corpus discusses but states no
value for --- and {{UNANS_OOD_REFUSED_PCT}} of the out-of-domain ones
({{UNANS_OOD_REFUSED_N}}), the latter being unanswerable for any system that
knows what the corpus is about.

% BEGIN INDUSTRIAL-SECTION
% The second-corpus section is generated from the industrial result files by
% deliverables/scripts/make_tex_industrial.py. It is prose, not a template:
% the sentences in it are the kind that go stale -- whether fusion beat its
% lexical leg, whether the band replicated, whether the rubric or the reader
% moved the number -- and each is a fact about the run rather than an
% assertion a writer may make either way.
% END INDUSTRIAL-SECTION

\section{Discussion}

The framework's strongest property is that each mechanism is \emph{structural}:
the gate is a branch rather than an instruction, the quota is enforced on the
artefact the generator consumes, and the conflict report is an orthogonal
conjunction rather than one classifier's opinion. Such properties are auditable,
and auditability is functional: a refusal whose reason can be reconstructed
from a score and a threshold is usable in a regulatory conversation. The second
strength is the explicit separation of coverage from correctness: a single
hallucination rate hides the trade-off. The third is the separation of
\emph{detection} from \emph{typing} in the contradiction gate: a single
precision number would have hidden the more consequential of the two, which
is that the mechanism finds real conflicts and then reports them under the
wrong label.

The framework does \emph{not} resolve contradictions; it surfaces them, which
the evidence suggests beats resolving. On the second corpus it goes further
and does not even say which version of a superseded provision is in force,
because nothing in the retrieved passages carries the version: a detector for
regulatory text needs version metadata, or must be scoped to detection alone.
It does not guarantee its own
confidence: the calibration is empirical, under an exchangeability assumption a
single-category corpus violates, so the conformal diagnostics are a limitation,
not a guarantee. The runtime form check does not parse English: being lexical,
it misses stems whose malformation is grammatical and cannot see a keyboard
mash. It refuses {{FORMVAL_STEM_MAL_RATE}} of the hand-labelled malformed stems
and fails open on what it cannot judge, the conservative direction.

{{GEN_LC_OPENING}} The most serious threat is a 2026 study reporting
long-context systems ahead of semantic RAG in manufacturing safety; our corpus
fits a modern context window, so the baseline is feasible. Experiment~7
measures the cost half: the whole corpus is {{LC_CORPUS_TOKENS}} tokens against
the RAG arm's {{LC_RAG_TOKENS}}, priced in dollars. {{GEN_LC_COST_CLAIM}}
The six-category corpus exceeds the 128K window most hosted models offer but
fits the generator's 1M-token window, so the answer-quality half is measured on
the focus domain, where the corpora are never pooled, not because the
comparison was impossible. {{GEN_LC_FINDING}}

Three results bear on this threat. Agri-Query~\cite{chen2026agriquery}, the
closest analogue to ours, favours retrieval: it places a 165-page
agricultural-machinery manual (roughly 59K tokens) in the prompt and reports
hybrid RAG beating direct long context, naive insertion significantly worse.
The 2025--2026 consensus is not \emph{which} to use but \emph{when to route}:
LaRA~\cite{lu2025lara} escalates only the queries that need it, and our
long-context arm is the always-long-context arm of that decision, measured
rather than assumed. Finally, \emph{Length Is Cheap, Competition Is Not}~\cite{liu2026lengthcheap}
predicts our own failure mode: attention mass on evidence explains about 58\%
of the variance in retrieval loss against at most 8\% for log-length, and one
competing value multiplies the loss by 9--110x, so a whole-corpus prompt fails
through competition among similar values, not length.
{{GEN_SUMMARY_SHORT}}

\section{Conclusion}

We presented an evidence-grounded RAG framework for industrial decision
support and were deliberate about each claim's strength. The architectural
contribution is incremental, positioned against published prior art rather than
asserted as priority. The empirical contribution is a failure-mode
characterisation: on an independently constructed public industrial corpus,
cross-encoder reranking is shown to improve recall only when its candidate
pool is an order of magnitude larger than the final context, the modality quota
to bind only before the reranker, the category pre-filter to cost no recall while improving the ordering, its benefit coming
from restricting the candidate pool the reranker sees rather than from reordering the fused
ranking, so that the soft-prior alternative proposed as its repair does not reproduce it, and the corpus's modality composition to
determine whether a modality-aware mechanism can have any effect at all. The
methodological contribution is an evaluation protocol treating the confidence
gate as a selective-prediction policy with a risk--coverage curve and a
false-positive rate at a stated operating point, with always-flag and never-flag
baselines.

Its value depends on each mechanism being structural, on their trade-offs being
measured rather than assumed, and on every claim naming its supporting passage.
On a second corpus of versioned federal regulations those negative findings
replicate --- the lexical leg carried the pipeline and neither the fusion nor
the category pre-filter added anything measurable on top of it --- and the
candidate-pool curve stops being saturated, so that the displacement of
retrieved evidence out of the context becomes directly measurable. The
clearest new result there is the decomposition: detection precision
{{IND_RO4_DETECTION}} against readers' labels and
{{IND_RO4_CONSTRUCTION_DETECTION}} against construction labels, against typing
precision {{IND_RO4_TYPING}} and {{IND_RO4_CONSTRUCTION_TYPING}}. A
$2\times2$ factorial then shows that the wording of the annotation rubric,
not the reader's domain specialisation, moved that measured precision, which
is a warning about the measurement rather than a clean result about the
mechanism.

\section*{Acknowledgment}
The author thanks Dr. Sri Krishna Kumar for supervision, and acknowledges the
Department of Industrial and Systems Engineering, IIT Kharagpur.

\begin{thebibliography}{99}
\bibitem{lewis2020rag} P. Lewis, E. Perez, A. Piktus, et al., ``Retrieval-augmented generation for knowledge-intensive NLP tasks,'' in \emph{NeurIPS}, 2020, pp. 9459--9474.
\bibitem{karpukhin2020dpr} V. Karpukhin, B. Oguz, S. Min, et al., ``Dense passage retrieval for open-domain question answering,'' in \emph{EMNLP}, 2020, pp. 6769--6781.
\bibitem{asai2024selfrag} A. Asai, Z. Wu, Y. Wang, A. Sil, and H. Hajishirzi, ``Self-RAG: Learning to retrieve, generate, and critique through self-reflection,'' in \emph{ICLR}, 2024.
\bibitem{yan2024crag} S.-Q. Yan, J.-C. Gu, Y. Zhu, and Z.-H. Ling, ``Corrective retrieval augmented generation,'' 2024. [Online]. Available: arXiv:2401.15884
\bibitem{jiang2023flare} Z. Jiang, F. F. Xu, L. Gao, et al., ``Active retrieval augmented generation,'' in \emph{EMNLP}, 2023.
\bibitem{su2024dragin} W. Su, Y. Tang, Q. Ai, and Z. Wu, ``DRAGIN: Dynamic retrieval augmented generation based on the real-time information needs of large language models,'' in \emph{ACL}, 2024.
\bibitem{jeong2024adaptive} S. Jeong, J. Baek, S. Cho, and S. J. Hwang, ``Adaptive-RAG: Learning to adapt retrieval-augmented LLMs through question complexity,'' in \emph{NAACL}, 2024.
\bibitem{rajpurkar2018squad2} P. Rajpurkar, R. Jia, and P. Liang, ``Know what you don't know: Unanswerable questions for SQuAD,'' in \emph{ACL}, 2018, pp. 518--528.
\bibitem{kadavath2022know} S. Kadavath, T. Conerly, A. Askell, and T. Henighan, ``Language models (mostly) know what they know,'' 2022. [Online]. Available: arXiv:2207.05221
\bibitem{kuhn2023semantic} L. Kuhn, Y. Gal, and S. Farquhar, ``Semantic uncertainty: Linguistic invariances for uncertainty estimation in natural language generation,'' in \emph{ICLR}, 2023.
\bibitem{farquhar2024semantic} S. Farquhar, J. Kossen, L. Kuhn, and Y. Gal, ``Detecting hallucinations in large language models using semantic entropy,'' \emph{Nature}, vol. 630, no. 8015, pp. 625--630, 2024.
\bibitem{joren2024sufficient} H. Joren, J. Zhang, C.-S. Ferng, D.-C. Juan, and A. Taly, ``Sufficient context: A new lens on retrieval augmented generation systems,'' 2024. [Online]. Available: arXiv:2411.06037
\bibitem{geissler2026dependable} F. Gei{\ss}ler, F. Carella, L. Fieback, and J. Spiegelberg, ``Towards dependable retrieval-augmented generation using factual confidence prediction,'' 2026. [Online]. Available: arXiv:2605.05244
\bibitem{phillips2026entropy} E. Phillips, F. Gustafsson, S. Wu, A. Thakur, and D. A. Clifton, ``Entropy alone is insufficient for safe selective prediction in LLMs,'' 2026. [Online]. Available: arXiv:2603.21172
\bibitem{thakur2026declinecraft} A. Thakur, A. Saparov, and B. Ribeiro, ``DeclineCraft: A RAG abstention benchmark for jointly measuring combined retrieval and parametric knowledge,'' 2026. [Online]. Available: arXiv:2610.28340
\bibitem{faysse2025colpali} M. Faysse, H. Sibille, T. Wu, et al., ``ColPali: Efficient document retrieval with vision language models,'' in \emph{ICLR}, 2025.
\bibitem{suri2024visdom} M. Suri, P. Mathur, F. Dernoncourt, et al., ``VisDoM: Multi-document QA with visually rich elements using multimodal retrieval-augmented generation,'' 2024. [Online]. Available: arXiv:2412.10704
\bibitem{strich2026t2} J. Strich, E. K. Isgorur, M. Trescher, et al., ``T2-RAGBench: Text-and-table benchmark for evaluating retrieval-augmented generation,'' in \emph{EACL}, 2026.
\bibitem{decian2026dynamic} J. Decian, C. Njeh, and H. Nakouri, ``Dynamic RAG with hybrid retrieval: A novel approach to robust LLMs and hallucination prevention,'' \emph{Neurocomputing}, vol. 705, art. 134911, 2026. https://doi.org/10.1016/j.neucom.2026.134911
\bibitem{gao2023alce} T. Gao, H. Yen, J.-T. Yu, and D. Chen, ``Enabling large language models to generate text with citations,'' in \emph{EMNLP}, 2023.
\bibitem{chen2026principled} Y. Chen, et al., ``Principled context engineering for RAG: Statistical guarantees via conformal prediction,'' in \emph{ECIR}, 2026. [Online]. Available: arXiv:2511.17908
\bibitem{longcontext2026tokentax} ``The token tax of epistemic accuracy: Comparing RAG and long-context architectures for document-grounded generative AI applications,'' 2026. [Online]. Available: arXiv:2606.20898
\bibitem{ranking2026notdeciding} ``Ranking is not deciding: Benchmark shortcuts and operating-point collapse in query-time RAG poisoning defenses,'' 2026. [Online]. Available: arXiv:2610.33849
\bibitem{claimconformal2026} ``Claim-level conformal prediction for retrieval-augmented generation,'' 2026. [Online]. Available: arXiv:2610.34872
\bibitem{feng2025conformal} N. Feng, Y. Sui, S. Hou, et al., ``Response quality assessment for retrieval-augmented generation via conditional conformal factuality,'' in \emph{SIGIR}, 2025. [Online]. Available: arXiv:2506.20978
\bibitem{syncheck2024} D. Wu, J.-C. Gu, F. Yin, and N. Peng, ``Synchronous faithfulness monitoring for trustworthy retrieval-augmented generation,'' in \emph{EMNLP}, 2024, pp. 9390--9406. [Online]. Available: \url{https://doi.org/10.18653/v1/2024.emnlp-main.527}
\bibitem{es2024ragas} S. Es, J. James, L. Espinosa-Anke, and S. Schockaert, ``RAGAS: Automated evaluation of retrieval augmented generation,'' in \emph{EACL System Demonstrations}, 2024.
\bibitem{saadfalcon2024ares} J. Saad-Falcon, O. Khattab, C. Potts, et al., ``ARES: An automated evaluation framework for retrieval-augmented generation systems,'' in \emph{NAACL}, 2024.
\bibitem{min2023factscore} S. Min, K. Krishna, X. Lyu, et al., ``FActScore: Fine-grained atomic evaluation of factual precision in long form text generation,'' in \emph{EMNLP}, 2023.
\bibitem{williams2018mnli} A. Williams, N. Nangia, P. Bowman, et al., ``An extensive evaluation of the natural language inference task,'' in \emph{ACL}, 2018.
\bibitem{nie2020anli} Y. Nie, A. H. Zhou, R. Socher, and J. Gao, ``ANLI: Annotating adversarially with natural language inference,'' in \emph{ACL}, 2020.
\bibitem{yang2026contextconflict} X. Yang, Z. Liu, R. Li, et al., ``Large language models in resolving contextual knowledge conflicts,'' in \emph{EMNLP}, 2026, pp. 21142--21169.
\bibitem{liu2026conflictscore} S. Liu, A. Halfaker, D. Roth, et al., ``ConflictScore: Identifying and measuring how language models handle conflicting evidence,'' in \emph{EMNLP}, 2026, pp. 58825--58845.
\bibitem{cormack2009rrf} G. V. Cormack, C. L. A. Clarke, and S. B\"uttcher, ``Reciprocal rank fusion outperforms Condorcet and individual rank learning methods,'' in \emph{SIGIR}, 2009, pp. 758--759.
\bibitem{bruch2023fusion} S. Bruch, S. Gai, and A. Ingber, ``An analysis of fusion functions for hybrid retrieval,'' \emph{ACM Trans. Inf. Syst.}, vol. 41, no. 6, 2023.
\bibitem{robertson2004prf} S. Robertson, H. Zaragoza, and J. Taylor, ``Probabilistic relevance framework: BM25 and beyond,'' \emph{Found. Trends Inf. Retr.}, vol. 3, no. 4, pp. 333--389, 2004.
\bibitem{huang2024selfcorrect} J. Huang, X. Chen, S. R. Mishra, et al., ``Large language models cannot self-correct reasoning yet,'' in \emph{ICLR}, 2024.
\bibitem{goswami2025haltrag} S. Goswami and S. Kurra, ``HALT-RAG: A task-adaptable framework for hallucination detection with calibrated NLI ensembles and abstention,'' 2025. [Online]. Available: arXiv:2509.07475
\bibitem{zong2026icalm} Y. Zong, Y. Li, X. Long, L. Chang, T. Wu, and R. Hadfield, ``I-CALM: Incentivizing confidence-aware abstention for LLM selective answering,'' 2026. [Online]. Available: arXiv:2604.03904
\bibitem{zhang2026taps} Y. Zhang, S. Wang, L. Yang, et al., ``Strong yet controllable LLM abstention without learning to abstain,'' 2026. [Online]. Available: arXiv:2610.16318
\bibitem{peng2025uaeval4rag} T. Peng, S. Choubey, X. Xiong, and W. Wu, ``Unanswerability evaluation for retrieval augmented generation,'' in \emph{ACL}, 2025. [Online]. Available: \url{https://aclanthology.org/2025.acl-long.415/}
\bibitem{tikhomirov2026selection} A. Tikhomirov and V. Savchenko, ``The selection effect: Why compact LLM verifiers for emotion classification fail behind encoder confidence gates,'' in \emph{EMNLP}, 2026. [Online]. Available: \url{https://aclanthology.org/2026.emnlp-main.1049/}
\bibitem{du2026recognition} S. Du and L. Hu, ``Recognition--refusal misalignment in LLMs: Why models answer structurally unanswerable questions,'' in \emph{EMNLP}, 2026. [Online]. Available: \url{https://aclanthology.org/2026.emnlp-main.1078/}
\bibitem{cocchieri2026overcommit} E. Cocchieri, L. Ragazzi, M. Tagliavini, and M. Moro, ``LLMs (almost) never abstain under medical uncertainty,'' in \emph{ACL}, 2026. [Online]. Available: \url{https://aclanthology.org/2026.acl-long.1365/}
\bibitem{chen2026agriquery} ``Agri-Query: A case study on RAG vs. long-context LLMs for cross-lingual technical question answering,'' 2025. [Online]. Available: arXiv:2508.18093
\bibitem{lu2025lara} ``LaRA: Benchmarking retrieval-augmented generation and long-context LLMs --- no silver bullet for LC or RAG routing,'' 2025. [Online]. Available: arXiv:2502.09977
\bibitem{liu2026lengthcheap} ``Length is cheap, competition is not: Attention predicts context rot in retrieval,'' 2026. [Online]. Available: arXiv:2611.04169
\bibitem{rosario2026jev} S. Do Rosario and M. Atef, ``Jev-as-a-judge for RAG claim verification: Confidence-gated cascading between a non-generative verifier and a frontier LLM,'' 2026. [Online]. Available: arXiv:2610.20634
\bibitem{ravuri2026rcl} ``RCL: A retrieval-confidence layer for detecting insufficient context in enterprise retrieval-augmented code generation,'' 2026. [Online]. Available: arXiv:2609.11023
\bibitem{zheng2026citeguard} ``CiteGuard-RAG: A validation-centered AI system for evidence-grounded question answering,'' 2026. [Online]. Available: arXiv:2609.15830
\bibitem{tian2026hira} Y. Tian, X. Chen, and L. Queiroz, ``HIRA: A human-in-the-loop retrieval-augmented cascade for document classification in regulated industries,'' in \emph{CIKM}, 2026. [Online]. Available: arXiv:2608.21792
\bibitem{chen2026modalityutility} ``Modality relevance is not modality utility: Post-hoc selective modality escalation for cost-aware multimodal RAG,'' 2026. [Online]. Available: arXiv:2607.05438
\bibitem{sharma2026placebo} ``Does document structure help dense retrieval? A placebo-controlled ablation of four mechanisms across two corpora,'' 2026. [Online]. Available: arXiv:2610.10170

\bibitem{yang2024ehrsql} Y. Yang, S. Kim, S. Kim, G. Lee, S.-Y. Yun, and E. Choi, ``Towards unbiased evaluation of detecting unanswerable questions in EHRSQL,'' 2024. [Online]. Available: arXiv:2405.01588
\bibitem{mondal2026beforeanswering} J. J. Mondal, M. S. A. Apurba, M. Banik, M. M. Al Mahmud, and I. F. Shihab, ``Before answering: Evidence sufficiency under size-matched memory construction,'' 2026. [Online]. Available: arXiv:2609.32269
\bibitem{brunello2026folio} A. Brunello, C. Curaba, L. Geatti, M. Mignani, A. Montanari, and N. Saccomanno, ``Fixing FOLIO and MALLS: Verified annotations and an LLM-assisted framework to focus human relabeling,'' 2026. [Online]. Available: arXiv:2606.02837
\bibitem{gautam2023unanswerable} V. Gautam, M. Zhang, and D. Klakow, ``A lightweight method to generate unanswerable questions in English,'' in \emph{Findings of EMNLP}, 2023. [Online]. Available: arXiv:2310.19403
\bibitem{semnani2025wikicollide} S. J. Semnani, J. Burapacheep, A. Khatua, et al., ``Detecting corpus-level knowledge inconsistencies in Wikipedia with large language models,'' in \emph{EMNLP}, 2025. [Online]. Available: arXiv:2509.23233
\bibitem{won2026sievesage} J. Won, S. G. An, and J.-Y. Lee, ``Sieve and Sage: Efficient distraction filtering for reliable RALM abstention,'' in \emph{Findings of EMNLP}, 2026. [Online]. Available: arXiv:2609.35794
\bibitem{li2026relevance} Y. Li, S. Li, Z. Xu, et al., ``Relevance is not sufficiency: What actually closes the evidence gap in long-term memory QA,'' 2026. [Online]. Available: arXiv:2610.09348
\bibitem{setiawan2026grabrag} Y. A. Setiawan, ``Prompt-based abstention fails under misleading context: A controlled study of small frozen RAG models,'' in \emph{AACL-IJCNLP}, 2026. [Online]. Available: arXiv:2608.22228
\bibitem{ling2026promptartifact} Z. Ling, S. Liu, Y. Tang, et al., ``LLM abstention can be a prompt artifact, in addition to genuine uncertainty,'' in \emph{EMNLP}, 2026, pp. 12612--12633. [Online]. Available: \url{https://doi.org/10.18653/2026.emnlp-main.549}
\bibitem{singh2026beyondreranker} S. Singh, A. Reddy, and M. Chopra, ``Beyond the reranker: Do RAG retrieval enhancements help once a strong reranker is present?'' 2026. [Online]. Available: arXiv:2606.28367
Y. Zhou, J. Yin, F. Zhang, et al., ``Financial evidence crowding: Diagnosing and mitigating constraint-induced displacement in retrieval-augmented generation,'' 2026. [Online]. Available: arXiv:2609.35782
\bibitem{zhang2026certify} X. Zhang and W. Sun, ``Knowing when the evidence is enough: Requirement-directed search with Horn-closure certificates,'' 2026. [Online]. Available: arXiv:2610.34684
\bibitem{qiu2026surerag} J. Qiu, Z. Han, and C. Huang, ``SURE-RAG: Sufficiency and uncertainty-aware evidence verification for selective retrieval-augmented generation,'' 2026. [Online]. Available: arXiv:2605.03534
\bibitem{wei2026kontrast} F. Wei, T. Ehrhart, and R. Troncy, ``Detecting knowledge inconsistencies across text, tables, and knowledge graphs,'' 2026. [Online]. Available: arXiv:2607.25959
\bibitem{valjakka2026labeling} J. Valjakka, J. Kivim\"aki, J. Myll\"ari, and J. K. Nurminen, ``The labeling problem in hallucination detection benchmarks: An empirical evaluation,'' in \emph{NeurIPS Datasets and Benchmarks}, 2026. [Online]. Available: arXiv:2610.08026
\bibitem{joshi2026realdocbench} A. Joshi, J. Kim, G. Eggert, et al., ``RealDocBench: A benchmark for field-level QA and layout understanding on real-world regulated documents,'' 2026. [Online]. Available: arXiv:2606.07401
\bibitem{park2026induspecqa} J. Park, H. Jeon, Y. Lee, and M. Kim, ``Ontology-based knowledge graphs for industrial standards with scope-aware logical rule structuring,'' in \emph{EMNLP}, 2026, pp. 33282--33311. [Online]. Available: \url{https://doi.org/10.18653/2026.emnlp-main.1456}
\bibitem{mori2025lexical} L. Mori, et al., ``Assessing the performance gap between lexical and semantic models for information retrieval with formulaic legal language,'' 2025. [Online]. Available: arXiv:2506.12895
\bibitem{lotfi2024bsard} E. Lotfi, et al., ``Bilingual BSARD: Extending statutory article retrieval to Dutch,'' in \emph{RegNLP}, 2025. [Online]. Available: arXiv:2412.07462
\bibitem{akarsu2026bm25} M. Akarsu, et al., ``From BM25 to corrective RAG: Benchmarking retrieval strategies for text-and-table documents,'' 2026. [Online]. Available: arXiv:2604.01733
\bibitem{beta2026greekbar} E. Beta, et al., ``GreekBarRetrieval: A benchmark for Greek statutory retrieval,'' in \emph{NLLP}, 2026. [Online]. Available: arXiv:2608.18752
\bibitem{gajjar2026teleembed} P. Gajjar and V. K. Shah, ``TeleEmbedBench: A multi-corpus embedding benchmark for RAG in telecommunications,'' 2026. [Online]. Available: arXiv:2604.17778
\bibitem{cheng2024silicon} X. Cheng, et al., ``To err is human; to annotate, SILICON? Toward robust reproducibility in LLM annotation,'' 2024. [Online]. Available: arXiv:2412.14461
\bibitem{kim2026guidelines} K. W. Kim, et al., ``Refining and reusing annotation guidelines for LLM annotation,'' in \emph{ACL}, 2026. [Online]. Available: arXiv:2605.20809
\bibitem{akarajaradwong2026nitibench} P. Akarajaradwong, et al., ``NitiBench+: A temporal Thai legal question answering benchmark across 17 domains,'' 2026. [Online]. Available: arXiv:2611.03576
\bibitem{prior2026oldfriend} M. Prior, et al., ``Asking for an old friend: Diagnosing and mitigating temporal failure modes in LLM-based statutory question answering,'' in \emph{ICAIL}, 2026. [Online]. Available: arXiv:2605.23497
\bibitem{huang2026wronglaw} Y. Huang, et al., ``When do LLMs apply the wrong law? Diagnosing LLM failures in temporal legal reasoning,'' 2026. [Online]. Available: arXiv:2608.14610
\bibitem{yu2026detecting} Z. Yu, et al., ``Detecting is not resolving: The monitoring control gap in retrieval augmented LLMs,'' 2026. [Online]. Available: arXiv:2605.27157
\bibitem{cattan2025dragged} A. Cattan, et al., ``DRAGged into conflicts: Detecting and addressing conflicting sources in search-augmented LLMs,'' 2025. [Online]. Available: arXiv:2506.08500
\end{thebibliography}

\end{document}
